\chapter{Nuclear Binding by Coulomb Alone}\label{ch:nuclear-binding}
\index{nuclear binding}\index{RealNucleus}\index{strong force}\index{QCD}\index{alpha ladder}\index{deuteron}\index{Standard Model}

\begin{quote}
\itshape If the same Coulomb energetics that binds an atom also binds a nucleus, the strong force is a name for
what confinement already does.\\
\normalfont\hfill --- working note, 2026.
\end{quote}
\bigskip

Chapter~\ref{ch:nucleus} took the structural step: the RealQM machinery of charge densities held in mutual
Coulomb confinement, applied at the nuclear scale, gives a He-4 nucleus as a small analogue of a closed-shell
atom. This chapter makes the step quantitative and confronts it with the incumbent theory. The claim is stark:
\textbf{nuclear binding energies follow from Coulomb interactions alone, with no strong or weak force} --- the
nucleus read as protons and electrons bound exactly as atoms are, only smaller. We report what the model
computes, how it compares with experiment, and how it stands against quantum chromodynamics; and we are explicit
that the account is structural and speculative, not a finished theory of the nucleus.

\section{The dual-Coulomb model}\label{sec:dual-coulomb}

The picture is the chemistry picture at a smaller length scale. Protons carry the positive kernels; a set of
electrons, confined to the nuclear volume, provides the negative charge whose Coulomb attraction to two or more
protons binds them --- the electron plays between protons the role of the shared bonding pair in a covalent
molecule. Each charge density is a phase of the same RealQM continuum, non-overlapping, held against collapse by
its own kinetic (localisation) energy exactly as in the atomic problem. Nothing is added to the energy functional
of Chapter~\ref{ch:equation}: it is Coulomb attraction, Coulomb repulsion, and the localisation cost, evaluated
at a femtometre rather than an \aa ngstr\"om.

\section{Results: helium-4, the alpha ladder, the deuteron ratio}\label{sec:nuc-results}

The simplest non-trivial nucleus, He-4, comes out bound by about $29.1$~MeV against the observed $28.3$~MeV ---
$103\%$, with no parameter tuned. The $\alpha$-conjugate ladder (He-4, C-12, O-16, \dots), treated as stacked
He-4-like units, is reproduced at the $107$--$108\%$ level. And the ratio of the He-4 binding to the deuteron
binding, a pure number free of the overall energy scale, comes out $13.1$ against the experimental $12.7$. That a
single confinement model, carried over unchanged from chemistry, lands the absolute binding, the ladder, and the
scale-free ratio together is the substance of the claim.

\section{One framework, two scales}\label{sec:two-scales}

The same variational principle thus spans two regimes that mainstream physics assigns to different forces:
electromagnetism for the atom, the strong interaction for the nucleus. In RealQM they are one Coulomb continuum
at two scales, related by a single ratio of lengths --- the atomic Bohr radius to the nuclear radius --- which is
essentially the fine-structure constant reappearing as a size ratio (Chapter~\ref{ch:cosmology}). The
proton and the electron exchange \emph{roles} but not \emph{physics}: in the atom the heavy kernels are protons
and the light glue is electrons; in the nucleus the heavy confined objects are protons and the binding glue is
again electrons, now compressed to nuclear dimensions. This role-symmetry --- the ``unified Coulomb'' reading of
atom and nucleus --- is what lets one code, one functional, address both.

\section{Against quantum chromodynamics}\label{sec:vs-qcd}

The contrast with the incumbent theory is sharp and worth stating plainly. Quantum chromodynamics, the accepted
theory of the strong interaction since 1973, has after half a century \emph{not} delivered nuclear binding
energies from first principles without fitting: lattice QCD computes light nuclei only at unphysically large
quark masses and extrapolates, and chiral effective field theory carries low-energy constants fitted to nuclear
data. The binding of the deuteron, and the $\alpha$ ladder, are not parameter-free QCD predictions. The
Coulomb-confinement model delivers the absolute He-4 binding, the ladder, and the He-4/deuteron ratio from a
single length scale and no fitted force. Whatever its ultimate standing, the model does the one thing QCD has not:
it predicts nuclear binding numbers without adjustable strong-force parameters.

Against the Standard Model more broadly, the question the model poses is one of economy: the Standard Model
carries a large tally of independent parameters and a separate term for each interaction, where RealQM proposes a
single Coulomb term at two scales. Extension or point, one term or many --- the model is a wager that the second
is enough for the binding of matter.

\section{Status: structural and speculative}\label{sec:nuc-binding-status}

We keep the honest label of Chapter~\ref{ch:nucleus}: this is \emph{structural and speculative}. The model
addresses \emph{binding}, not \emph{substructure}: it treats protons and electrons as Poisson charge densities
and says nothing about the parton content probed in deep-inelastic scattering, which is a genuinely different,
high-energy regime outside its scope. It does not yet compute the full chart of nuclides in detail, and the
role-symmetric confinement of electrons at nuclear scale raises questions --- of stability, of the electron's
fate --- that a finished theory would have to answer. What is established is narrower and still remarkable: a
chemistry-derived Coulomb model, carried unchanged to the femtometre scale, reproduces the principal nuclear
binding numbers that the strong-force theory obtains only by fitting. That is a fact worth the scrutiny it
invites, offered as such.
